% !TeX root = ../main.tex

\begin{itemize}
	\item Suivre 2 axes : les contraintes vis-à-vis de la gestion des collections, donc comment concilier besoin utilisateur et contraintes techniques, temporelles et d'ingénierie documentaire. Et axe 2 sur une contrainte plus technique qui est le management de ressources informatiques sur des collections titanesques notamment dans le cadre du multimédia. 
	% \item  L’IA peut-elle faire face à une telle diversité de données ? 
	% \item  L’IA peut-elle être utilisée sur toutes les collections ? Mêmes les plus grandes et diversifiées ? 
	% \item  Peut-on penser intégrer de l’IA tout en conservant et respectant l’expérience utilisateur actuelle ?
	
	\item \textbf{Plan :}
	\subitem 1. L'IA a une capacité de traitement énorme, mais qui a un coût financier et écologique, faire un point sur le sujet. 
	\subitem 2. Insister sur les enjeux d'infrastructures : serveurs, RAM, VRAM, CPU. Parfois les institutions n'ont pas la possibilité d'accueillir de telles infrastructures (contraire à la politique mise en place par le SI, trop chère, espace, etc.) 
	\subitem 3. Evolutions extrêmement rapides de l'IA dans tous les domaines qu'elle regroupe, cela signifie aussi une augmentation constante des prix pour l'accueillir. Utiliser une IA quel qu'elle soit c'est un calcul financier, technique et politique. 
	\subitem 4. Dans le monde patrimonial la taille des collections amène des traitements monumentaux. L'enjeu financier va certainement amener à rationaliser les attentes, établir des priorités de traitement et donc morceler ces collections.
	\subitem 5. Intégrer l'IA c'est aussi prendre le risque de modifier un chemin utilisateur, une habitude utilisateur ou d'aller contre une volonté utilisateur. En plus d'être pensée selon son coût et besoin technique, l'implémentation de l'IA doit aussi être réfléchie sous l'angle de la réception utilisateur, de l'impact sur la qualité d'usage de l'applicatif, etc.
	
\end{itemize}